\section{Every logarithmic moment at arbitrary complex orders}
\label{rlc:sec:moments}

Order addition alone evaluates the unweighted convolution. Polynomial
weights record how the total logarithmic argument is distributed among
its factors. Surprisingly, they still admit finite closure at arbitrary
complex orders.

\subsection{An all-depth differential-polynomial algorithm}
For $\ell\ge0$ define polynomials $R_\ell(s;T,U)$ by
\begin{align}
 R_0&=1,\notag\\
 R_{\ell+1}&=(T+sU)R_\ell+(\pi^2+T^2)\partial_T R_\ell
                        +U^2\partial_U R_\ell.
 \label{rlc:eq:Rrecurrence}
\end{align}
Their first values are
\begin{align*}
 R_1&=T+sU,\\
 R_2&=2T^2+2sTU+s(s+1)U^2+\pi^2,\\
 R_3&=6T^3+6sT^2U+3s(s+1)TU^2+(s)_3U^3
          +5\pi^2T+3s\pi^2U.
\end{align*}
For a multi-index $\boldsymbol\ell=(\ell_1,\ldots,\ell_r)$, let
\begin{equation}\label{rlc:eq:weighteddef}
 \MM_{\boldsymbol\ell}(\boldsymbol s;a;\mu)
 =\int_{\R^{r-1}}\prod_{j=1}^r
     x_j^{\ell_j}f_{a,s_j}(x_j)\dd x_1\cdots\dd x_{r-1},
 \qquad x_r=\mu-\sum_{j<r}x_j.
\end{equation}

\begin{theorem}[Finite closure of all mixed moments]
\label{rlc:thm:allmoments}
The integral \eqref{rlc:eq:weighteddef} is absolutely convergent and
entire in its orders. Form the finite polynomial
$\prod_{j=1}^r R_{\ell_j}(s_j;T,U)$ and put $W=\sum s_j$.
Apply the following linear replacement rules to its monomials:
\begin{align}
 T^{2k}U^d\ &\longmapsto\
 \sum_{q=0}^k\binom{k}{q}(-\pi^2)^{k-q}
        \CC_{r+2q,a}(W+d;\mu),\label{rlc:eq:evenrule}\\
 T^{2k+1}U^d\ &\longmapsto\
 \sum_{q=0}^k\binom{k}{q}\frac{(-\pi^2)^{k-q}}{r+2q}
 \left\{\mu\CC_{r+2q,a}(W+d;\mu)
 -(W+d)\CC_{r+2q,a}(W+d+1;\mu)\right\}.
 \label{rlc:eq:oddrule}
\end{align}
The resulting expression is exactly
$\MM_{\boldsymbol\ell}(\boldsymbol s;a;\mu)$.
Thus every mixed logarithmic moment at $\mu=0$ reduces to finitely many
entire Hurwitz-zeta combinations in even depth, and alternating
Hurwitz-zeta combinations in odd depth. Arbitrary mixed order
derivatives of this finite formula are valid.
\end{theorem}
\begin{proof}
Put $g(p)=\pi\csc(\pi p)$, $h=p+a-1$, $T=\pi\cot(\pi p)$,
and $U=h^{-1}$. Then
$g'=-Tg$, $T'=-(\pi^2+T^2)$, and $U'=-U^2$.
Repeated differentiation gives the exact identity
\[
 (-\partial_p)^\ell\{g(p)h^{-s}\}
       =g(p)h^{-s}R_\ell(s;T,U).
\]
The left side is the transform of $x^\ell f_{a,s}(x)$.
The product of the transforms of all factors is consequently
$g(p)^r h^{-W}\prod_jR_{\ell_j}(s_j;T,U)$.

The identity $T^2=g^2-\pi^2$ proves
\eqref{rlc:eq:evenrule}. For an odd power, first use that identity
and then $g^R T=-(g^R)'/R$. Integration by parts on the vertical line
in \eqref{rlc:eq:barnes} gives, for every $v$,
\[
 \frac1{2\pi i}\int e^{p\mu}g(p)^R T h^{-v}\dd p
 =\frac{\mu}{R}\CC_{R,a}(v;\mu)
       -\frac{v}{R}\CC_{R,a}(v+1;\mu).
\]
The boundary terms vanish exponentially. This proves the odd rule.
Lemma~\ref{rlc:lem:weighted} supplies absolute convergence and parameter
holomorphy, so no formal interchange remains unproved.
\end{proof}

For example, at $\mu=0$, with $s=s_1$ and $W=\sum s_j$,
\begin{align}
 \MM_{(1,0,\ldots,0)}
 &=\left(s-\frac Wr\right)\CC_{r,a}(W+1;0),
 \label{rlc:eq:firstall}\\
 \MM_{(2,0,\ldots,0)}
 &=2\CC_{r+2,a}(W;0)-\pi^2\CC_{r,a}(W;0)\notag\\
 &\quad+\left[s(s+1)-\frac{2s(W+1)}r\right]
                \CC_{r,a}(W+2;0).
 \label{rlc:eq:secondall}
\end{align}
For two distinct positions having orders $s$ and $t$,
\begin{align}
 \MM_{(1,1,0,\ldots,0)}
 &=\CC_{r+2,a}(W;0)-\pi^2\CC_{r,a}(W;0)\notag\\
 &\quad+\left[st-\frac{(s+t)(W+1)}r\right]
                    \CC_{r,a}(W+2;0).
 \label{rlc:eq:mixedsecond}
\end{align}
Summing \eqref{rlc:eq:firstall} over all positions gives zero, as it
must from $x_1+\cdots+x_r=0$. The factor $1/r$ is essential.

\subsection{A closed two-factor formula}
For $k\ge1$ define
\begin{equation}\label{rlc:eq:Q}
 Q_k(s,t)=\frac{(s)_k-(-1)^k(t)_k}{s+t+k-1}.
\end{equation}
This is a polynomial, not a meromorphic quotient: at
$t=1-k-s$, $(t)_k=(-1)^k(s)_k$, proving divisibility.
Its first values are
\[
 Q_1=1,\qquad Q_2=s-t,\qquad
 Q_3=s^2-st+t^2+s+t.
\]
Let $d_{2j}$ be defined by
\begin{equation}\label{rlc:eq:dcoeff}
 \frac{\pi\lambda}{\sin\pi\lambda}
 =\sum_{j=0}^\infty d_{2j}\lambda^{2j},\qquad
 d_0=1,\quad
 d_{2j}=2(1-2^{1-2j})\zeta(2j)\quad(j\ge1).
\end{equation}

\begin{theorem}[All two-factor logarithmic moments]
\label{rlc:thm:twomoments}
For $\ell\ge0$, $s,t\in\C$, and $\Rea a>0$, put $W=s+t$. Then
\begin{equation}\label{rlc:eq:momentclosed}
 \boxed{\displaystyle
 \int_0^\infty\log^\ell y\,F_{a,s}(y)F_{a,t}(1/y)\frac{\dd y}{y}
 =\ell!\sum_{j=0}^{\lfloor\ell/2\rfloor}
 \frac{d_{2j}}{(\ell-2j+1)!}
 Q_{\ell-2j+1}(s,t)\E_a(W+\ell-2j).}
\end{equation}
Every term on the right is entire in the two independent orders.
\end{theorem}
\begin{proof}
Introduce the locally holomorphic tilted integral
\[
 I(\lambda)=\int_0^\infty y^\lambda
 F_{a,s}(y)F_{a,t}(1/y)\frac{\dd y}{y}.
\]
The weighted estimates prove convergence for
$|\Rea\lambda|<\min(1,\Rea a)$. For
$\Rea s,\Rea t>0$ and sufficiently small $|\lambda|$, inserting
Euler integrals and evaluating the elementary rational $y$ integral
gives
\begin{equation}\label{rlc:eq:tilted}
 I(\lambda)=\frac{\pi}{\sin\pi\lambda}
 \sum_{n=0}^\infty\left[
 (n+a-\lambda)^{-s}(n+a)^{-t}
 -(n+a)^{-s}(n+a+\lambda)^{-t}\right].
\end{equation}
For clarity, the elementary integration kernel before the two Laplace
integrals are performed is
\[
 \frac{\pi}{\sin\pi\lambda}
 \frac{e^{\lambda u}-e^{-\lambda v}}{e^{u+v}-1}
 e^{-(a-1)(u+v)}.
\]
It is multiplied by $u^{s-1}v^{t-1}/(\Gamma(s)\Gamma(t))$.
Expanding the geometric denominator and integrating $u,v$ gives
\eqref{rlc:eq:tilted}. Its $\lambda=0$ value is removable.

Taylor expansion in $\lambda$ now yields
\[
 I(\lambda)=\frac{\pi}{\sin\pi\lambda}
 \sum_{k=1}^\infty
 \frac{(s)_k-(-1)^k(t)_k}{k!}
       \zeta(W+k,a)\lambda^k.
\]
Normal convergence in a small disk follows from $\Rea W>0$ and the
shift bound $|\lambda|<\Rea a$. Extracting the coefficient of
$\lambda^\ell$ and using
\[
 \bigl((s)_k-(-1)^k(t)_k\bigr)\zeta(W+k,a)
       =Q_k(s,t)\E_a(W+k-1)
\]
proves \eqref{rlc:eq:momentclosed} in the initial domain. Both sides
are entire in $s,t$, which proves the theorem everywhere.
\end{proof}

The first two nonconstant moments are therefore
\begin{align}
 \int_0^\infty\log y\,F_{a,s}(y)F_{a,t}(1/y)\frac{\dd y}{y}
 &=\frac{s-t}{2}\E_a(W+1),\label{rlc:eq:moment1}\\
 \int_0^\infty\log^2 y\,F_{a,s}(y)F_{a,t}(1/y)\frac{\dd y}{y}
 &=\frac{s^2-st+t^2+s+t}{3}\E_a(W+2)
      +\frac{\pi^2}{3}\E_a(W).
 \label{rlc:eq:moment2}
\end{align}
In particular, at $s=t=0$ the second moment is $\pi^2/3$, the ordinary
logistic second moment. Formula \eqref{rlc:eq:moment1} independently
checks the $1/2$ in the first-moment normalization. All odd moments
vanish when $s=t$, directly reflecting $y\leftrightarrow1/y$.
